\begin{table}[htbp]
\centering
\small
\caption{Impedance reference distances.}
\label{tab:t4_reference}
\begin{tabular}{lll}
\toprule
Mode & $T_m$ & Basis \\
\midrule
Major road & 2 km & UDPFI/URDPFI: arterial minimum intersection spacing 500 m; sub-arterial spacing 0.5 km (CBD) to 3–5 km (residential). Half the mean spacing implies a catchment edge of 1–2.5 km. \\
ICD & 50 km & CBIC Circular No. 50/2020-Customs, para 4.1.4: no greenfield ICD within 100 km of an existing ICD. Statutory 100 km spacing implies a ~50 km catchment under nearest-facility assignment. \\
Rail goods shed & 50 km & Weakest anchor. Intermodal first-mile drayage is conventionally 50–100 km. No Indian Railways or CONCOR catchment norm identified; stated as an assumption. \\
IWT terminal & 50 km & By analogy with ICD. No statutory anchor: the CBIC circular relaxes distance rules for waterway-linked facilities. \\
\bottomrule
\end{tabular}
\begin{tablenotes}\footnotesize
\item $T_m$ is the catchment edge: the distance at which access is judged half that of co-location. All four are swept in the sensitivity analysis (Table S2).
\end{tablenotes}
\end{table}
